% Generated by scripts/build_latex.py from docs/report/appendix_a_user_manual.md: edit the Markdown, not this file.
\chapter{User Manual (Short)}\label{ch:A}

This appendix summarises how to install and use IDSS. The full manual is \code{docs/\allowbreak{}user\_\allowbreak{}manual.\allowbreak{}md} in the repository.

\FloatBarrier
\section{Requirements}\label{sec:A.1}

\begin{itemize}
  \item Windows 10 or 11, Python 3.12 and Git; about 2 GB of free disk space for the virtual environment.
  \item Optional: an API key for an OpenAI-compatible LLM provider (OpenAI, Groq, Google Gemini) or a local Ollama model. Without one, everything except \textit{Ask Your Data} works in offline mode.
\end{itemize}

\FloatBarrier
\section{Installation}\label{sec:A.2}

From a PowerShell window in the project folder:

\begin{lstlisting}
python -m venv .venv
.venv\Scripts\python.exe -m pip install -r requirements.txt
Copy-Item .env.example .env
# optional check: the offline test suite
.venv\Scripts\python.exe -m pytest -q
\end{lstlisting}

To use an LLM, set \code{LLM\_\allowbreak{}PROVIDER}, \code{LLM\_\allowbreak{}MODEL} and \code{LLM\_\allowbreak{}API\_\allowbreak{}KEY} in \code{.env} (the file stays on the computer and is ignored by Git). Leave them empty for offline mode.

\FloatBarrier
\section{Starting and Stopping}\label{sec:A.3}

\begin{lstlisting}
powershell -ExecutionPolicy Bypass -File scripts\run_all.ps1
\end{lstlisting}

This starts the API (\code{http:/\allowbreak{}/\allowbreak{}localhost:8000}, interactive documentation at \code{/docs}) in a new window and then the user interface at \code{http:/\allowbreak{}/\allowbreak{}localhost:8501}. Stop the UI with Ctrl+C and close the API window to stop the backend. Both servers listen on this computer only.

\FloatBarrier
\section{Pages}\label{sec:A.4}

\begin{table}[H]
\centering
\linespread{1}\small
\setlength{\tabcolsep}{4pt}
\caption{Pages of the user interface}\label{tab:A.1}
\begin{tabular}{|>{\raggedright\arraybackslash\hspace{0pt}}p{\dimexpr 0.2128\linewidth-2\tabcolsep-2\arrayrulewidth\relax}|>{\raggedright\arraybackslash\hspace{0pt}}p{\dimexpr 0.7872\linewidth-2\tabcolsep-2\arrayrulewidth\relax}|}
\hline
\rowcolor{tableheader}\textbf{Page} & \textbf{What the user does} \\
\hline
Home & Reads the overview, sees headline numbers and recent runs; \textbf{Demo Mode} runs the churn example in one click \\ \hline
Upload Data & Uploads a CSV, Excel, JSON or Parquet file (\ensuremath{\leq} 50 MB) and reads its profile and warnings \\ \hline
New Analysis & Describes the problem in plain language (or picks an example), chooses the \textit{fast} or \textit{full} budget and, optionally, the actionable columns, and starts the analysis \\ \hline
Run Monitor & Follows the stages and the event timeline; answers clarifying questions when the run asks for them \\ \hline
Results & Reads the test metrics, the leaderboard, the verification checks, the key drivers and the charts \\ \hline
Decision Center & Reads the decision brief, risk tiers or forecast and recommended actions; runs what-if simulations for one record or a segment \\ \hline
Ask Your Data & Asks questions about the data (needs an LLM); the generated code and its output are shown \\ \hline
Reports & Downloads the PDF or Markdown report and previews it \\ \hline
\end{tabular}
\end{table}

\FloatBarrier
\section{A First Run (about 3 minutes)}\label{sec:A.5}

\begin{enumerate}
  \item On \textbf{Home}, press \textbf{Demo Mode} (or upload \code{data/\allowbreak{}samples/\allowbreak{}customer\_\allowbreak{}churn.\allowbreak{}csv} on \textbf{Upload Data} and choose the churn example on \textbf{New Analysis}).
  \item On \textbf{Run Monitor}, watch the run move from INIT through TRAIN and VERIFY to DONE.
  \item On \textbf{Results}, check the best model, its F1 score on the held-out test data and the verification checks.
  \item In the \textbf{Decision Center}, read the High / Medium / Low risk tiers and the recommended actions, then open the what-if simulator, change \code{contract\_\allowbreak{}type} to another contract and press \textbf{Simulate}.
  \item On \textbf{Reports}, download the PDF.
\end{enumerate}

\FloatBarrier
\section{Troubleshooting}\label{sec:A.6}

\begin{table}[H]
\centering
\linespread{1}\small
\setlength{\tabcolsep}{4pt}
\caption{Common problems}\label{tab:A.2}
\begin{tabular}{|>{\raggedright\arraybackslash\hspace{0pt}}p{\dimexpr 0.4362\linewidth-2\tabcolsep-2\arrayrulewidth\relax}|>{\raggedright\arraybackslash\hspace{0pt}}p{\dimexpr 0.5638\linewidth-2\tabcolsep-2\arrayrulewidth\relax}|}
\hline
\rowcolor{tableheader}\textbf{Problem} & \textbf{Solution} \\
\hline
``Backend not running -- start scripts/run\_api.ps1'' & Start the API (or use \code{run\_\allowbreak{}all.ps1}) and refresh the page \\ \hline
Upload rejected (413 / 415) & Files must be at most 50 MB and of type \code{.csv}, \code{.xlsx}, \code{.json} or \code{.parquet} \\ \hline
``409 busy'' / ``Another run is training'' & Only one run trains at a time; wait for it on the Run Monitor page \\ \hline
Run stuck in NEEDS\_INPUT & It is waiting for an answer on the Run Monitor page \\ \hline
Results say ``not verified'' & A quality check failed (see Results \ensuremath{\rightarrow} Verification); treat the results with caution, try the \textit{full} budget or more data \\ \hline
Ask Your Data is disabled & Configure an LLM in \code{.env} and restart the API \\ \hline
A run is marked FAILED after a restart & The API stopped while it was training; start a new run \\ \hline
\end{tabular}
\end{table}
